\addchap{Notation}

The chapters share the notation below. A chapter that needs a clashing local symbol says so where it introduces it, and the notes at the end of each chapter translate the notation of its main sources into this one.

\addsec*{Data, models and losses}

\begin{tabular}{@{}p{0.27\textwidth}p{0.69\textwidth}@{}}
$X^n=(X_1,\dots,X_n)$ & the data, independent with distribution $q$, the truth \\
$(x_i,y_i)$ & an input and an output in a regression $y=f(w,x)+\varepsilon$, with $\varepsilon\sim\Normal(0,\sigma^2)$ and $\sigma=1$ unless stated \\
$p(x\mid w)$, $w\in W\subset\R^d$ & a model and its parameter; $W$ is compact unless stated \\
$\varphi(w)$ & the prior density, uniform on $[-1,1]^d$ unless stated \\
$a,b$; $a_1,\dots,a_L$; $A,B$ & the parameters of the running models (two scalars; a chain; two matrices) \\
$L_n(w)$, $L(w)$ & the empirical loss $-\frac1n\sum_i\log p(X_i\mid w)$ and its expectation \\
$S$, $S_n$ & the entropy of $q$ and its empirical version $-\frac1n\sum_i\log q(X_i)$ \\
$w_0$, $W_0$ & a true (optimal) parameter, and the set of all of them \\
$K(w)$, $K_n(w)$ & $L(w)-L(w_0)\ge0$, the Kullback--Leibler divergence $\KL(q\,\|\,p(\cdot\mid w))$ in the realizable case, and $L_n(w)-L_n(w_0)$
\end{tabular}

\addsec*{Posteriors and free energies}

\begin{tabular}{@{}p{0.27\textwidth}p{0.69\textwidth}@{}}
$\beta$ & an inverse temperature; the tempered posterior has density proportional to $\varphi(w)e^{-n\beta L_n(w)}$ \\
$\Ew{f}$, $\Ewb{f}$ & the posterior expectation of $f(w)$ at $\beta=1$ and at $\beta$ \\
$p^*(x)$ & the Bayes predictive density $\Ew{p(x\mid w)}$ \\
$F_n$, $F_n(\beta)$ & the free energy $-\log\int\varphi(w)e^{-nL_n(w)}\dd w$ (the negative log marginal likelihood), and its tempered version \\
$F_n-nS_n$ & the normalized free energy $-\log\int\varphi(w)e^{-nK_n(w)}\dd w$ \\
$\bar F(n)$ & the population free energy $-\log\int\varphi(w)e^{-nK(w)}\dd w$
\end{tabular}

\addsec*{The learning coefficient}

\begin{tabular}{@{}p{0.27\textwidth}p{0.69\textwidth}@{}}
$\zeta(z)$ & the zeta function $\int K(w)^z\varphi(w)\dd w$ \\
$\lambda$, $m$ & the learning coefficient (the real log canonical threshold of $K$ with respect to $\varphi$) and its multiplicity: $-\lambda$ is the largest pole of $\zeta$ and $m$ its order \\
$V(\eps)$, $v(t)$ & the volume $\int_{K(w)<\eps}\varphi(w)\dd w$ and the state density $V'(t)$ \\
$\lambda(w^*)$, $\hat\lambda(w^*)$ & the local learning coefficient at a point $w^*$ and an estimate of it \\
$\nu$ & the singular fluctuation
\end{tabular}

\addsec*{Generalization and criteria}

\begin{tabular}{@{}p{0.27\textwidth}p{0.69\textwidth}@{}}
$G_n$, $T_n$ & the Bayes generalization loss $-\E_X[\log p^*(X)]$ and the Bayes training loss $-\frac1n\sum_i\log p^*(X_i)$; $G_n-S$ is the generalization error \\
$V_n$ & the functional variance $\sum_i\Var_w[\log p(X_i\mid w)]$ \\
$\WAIC$, $\WBIC$ & $T_n+V_n/n$, and $\Ewb{nL_n(w)}$ at $\beta=1/\log n$; both on the scale of the loss, never multiplied by $2n$ \\
$\AIC$, $\BIC$ & the classical criteria, written on the same scales
\end{tabular}

\addsec*{Probability and asymptotics}

$\Prob$ and $\E$ without a subscript are over the data (and fresh test points); a subscript names other randomness. $\Normal(\mu,\sigma^2)$ is the normal distribution and $\Unif$ the uniform one. $O_p$ and $o_p$ refer to $n\to\infty$ with the randomness of the data; $a_n\sim b_n$ means $a_n/b_n\to1$, and $a\asymp b$ means that each is bounded by a constant multiple of the other. Logarithms are natural. $\N$ is the set of natural numbers.
